% appendices/appendix_a_symbols.tex

\chapter{数学符号与张量维度约定}

在大模型工程中，许多问题表面上看是系统问题，实际根源却来自张量维度、矩阵乘法方向、并行切分方式或内存布局理解不一致。例如，同样是一个 attention tensor，在论文中可能写作 $[B, H, S, D_h]$，在框架实现中可能变成 $[S, B, H, D_h]$，在 CUDA kernel 中又可能被重排为适合访存合并的连续布局。本附录统一本书使用的数学符号、维度命名和公式排版规则，方便读者在阅读后续章节时快速对齐语义。

\section{常用数学符号}

\subsection{标量、向量、矩阵与张量}

本书使用如下约定：

\begin{longtable}{p{0.22\textwidth}p{0.68\textwidth}}
\toprule
符号 & 含义 \
\midrule
$a, b, c$ & 标量，通常表示单个数值，例如 loss、学习率、温度参数等。 \
$\vect{x}, \vect{y}, \vect{z}$ & 向量，通常表示一个 token 的 hidden state、一个参数向量或一个概率分布。 \
$\mat{A}, \mat{B}, \mat{W}$ & 矩阵，通常表示线性层权重、二维激活或批量矩阵乘法中的单个矩阵。 \
$\tensor{X}, \tensor{Y}, \tensor{Z}$ & 高维张量，通常表示 batch、sequence、head、hidden 等多个维度组成的数据。 \
$\theta$ & 模型参数集合。 \
$\nabla_{\theta}\mathcal{L}$ & loss 对模型参数的梯度。 \
$\mathcal{L}$ & 损失函数。 \
$\eta$ & 学习率。 \
\bottomrule
\end{longtable}

在具体实现中，一个变量到底是向量、矩阵还是高维张量，取决于上下文。例如，Transformer 中的 hidden states 在单个 token 视角下是向量，在单个序列视角下是矩阵，在 batch 视角下是三维张量。

\subsection{索引与集合}

本书常用如下索引：

\begin{longtable}{p{0.22\textwidth}p{0.68\textwidth}}
\toprule
符号 & 含义 \
\midrule
$i, j, k$ & 通用索引，常用于矩阵元素、token 位置或循环变量。 \
$b$ & batch 中的样本索引。 \
$s, t$ & sequence 中的 token 位置索引。 \
$h$ & attention head 索引。 \
$l$ & Transformer 层索引。 \
$r$ & 分布式训练中的 rank 索引。 \
$N$ & 通用数量，例如样本数、节点数、GPU 数等。 \
$V$ & 词表大小 vocabulary size。 \
\bottomrule
\end{longtable}

若存在歧义，本书会在公式前显式说明索引含义。例如：
[
\tensor{X}_{b,s,:}
]
表示第 $b$ 个样本、第 $s$ 个 token 对应的 hidden vector。

\section{Transformer 中的维度命名}

Transformer 工程实现中最容易出错的地方之一，是不同维度含义相近但用途不同。尤其在多头注意力、KV Cache、张量并行和推理服务章节中，维度命名必须保持一致。

\subsection{基础维度}

本书默认使用如下维度符号：

\begin{longtable}{p{0.2\textwidth}p{0.24\textwidth}p{0.46\textwidth}}
\toprule
符号 & 名称 & 含义 \
\midrule
$B$ & batch size & 一个 step 或一次推理迭代中同时处理的样本数量。 \
$S$ & sequence length & 输入序列长度，也可表示 prefill 阶段上下文长度。 \
$T$ & generated length & 自回归生成阶段生成的 token 数。 \
$D$ & hidden size & 每个 token 的隐藏状态维度。 \
$H$ & number of heads & attention head 数量。 \
$D_h$ & head dimension & 单个 attention head 的维度，通常满足 $D = H \cdot D_h$。 \
$D_{\mathrm{ff}}$ & FFN intermediate size & Feed-Forward Network 中间层维度。 \
$V$ & vocabulary size & tokenizer 对应的词表大小。 \
$L$ & number of layers & Transformer block 层数。 \
\bottomrule
\end{longtable}

典型 Decoder-only LLM 的主要张量形状如下：

\begin{longtable}{p{0.28\textwidth}p{0.28\textwidth}p{0.34\textwidth}}
\toprule
对象 & 形状 & 说明 \
\midrule
Token IDs & $[B, S]$ & 输入 token 编号。 \
Embedding Output & $[B, S, D]$ & token embedding 后的隐藏状态。 \
Q / K / V & $[B, S, H, D_h]$ & 多头注意力中的 query、key、value。 \
Attention Scores & $[B, H, S, S]$ & prefill 阶段完整 self-attention 分数矩阵。 \
Attention Output & $[B, S, H, D_h]$ & 每个 head 的 attention 输出。 \
Merged Attention Output & $[B, S, D]$ & 多头合并后的输出。 \
FFN Intermediate & $[B, S, D_{\mathrm{ff}}]$ & FFN 上投影后的中间激活。 \
Logits & $[B, S, V]$ & 每个 token 位置上的词表 logits。 \
\bottomrule
\end{longtable}

\subsection{MHA、MQA 与 GQA 中的 head 维度}

标准 Multi-Head Attention 中，Query、Key、Value 通常具有相同数量的 heads：
[
H_q = H_k = H_v = H.
]

在 Multi-Query Attention 中，多个 query heads 共享一组 key/value heads：
[
H_q = H,\qquad H_k = H_v = 1.
]

在 Grouped-Query Attention 中，query heads 被分组，每组共享一个 key/value head：
[
H_q = H,\qquad H_k = H_v = H_{\mathrm{kv}},\qquad H_{\mathrm{kv}} < H.
]

因此，KV Cache 的大小与 $H_{\mathrm{kv}}$ 直接相关，而不是只与 $H_q$ 相关。推理系统中估算显存时，应使用：
[
\mathrm{KVCacheBytes}
=====================

2 \cdot B \cdot S \cdot L \cdot H_{\mathrm{kv}} \cdot D_h \cdot \mathrm{bytes_per_element}.
]
其中系数 $2$ 分别对应 Key Cache 与 Value Cache。

\section{batch、sequence、hidden、head 维度约定}

\subsection{默认逻辑布局}

除非特别说明，本书在数学推导中默认使用 batch-first 布局：
[
\tensor{X} \in \mathbb{R}^{B \times S \times D}.
]

这与许多现代 PyTorch 模型实现保持一致，便于读者理解 batch、sequence 与 hidden 的关系。在讲解注意力时，本书通常将 Q/K/V 写作：
[
\tensor{Q}, \tensor{K}, \tensor{V}
\in
\mathbb{R}^{B \times S \times H \times D_h}.
]

实际工程中，为了适配 kernel 或提高访存效率，常见布局还包括：

\begin{itemize}
\item $[S, B, D]$：早期序列模型和部分框架内部实现常用。
\item $[B, H, S, D_h]$：attention kernel 中常用，方便按 head 并行。
\item $[B, S, H \cdot D_h]$：linear projection 前后的合并 hidden 表示。
\item $[L, B, H_{\mathrm{kv}}, S, D_h]$：KV Cache 的一种逻辑表达方式。
\end{itemize}

本书会在涉及 kernel 或内存布局时明确区分“逻辑形状”和“物理布局”。逻辑形状描述语义，物理布局描述内存中连续存储的顺序。

\subsection{矩阵乘法方向}

线性层通常写作：
[
\mat{Y} = \mat{X}\mat{W} + \vect{b},
]
其中：
[
\mat{X} \in \mathbb{R}^{N \times D_{\mathrm{in}}},
\qquad
\mat{W} \in \mathbb{R}^{D_{\mathrm{in}} \times D_{\mathrm{out}}},
\qquad
\mat{Y} \in \mathbb{R}^{N \times D_{\mathrm{out}}}.
]

在 PyTorch 的 \texttt{nn.Linear(in_features, out_features)} 中，权重参数的存储形状通常是：
[
[D_{\mathrm{out}}, D_{\mathrm{in}}],
]
前向计算在语义上等价于：
[
\mat{Y} = \mat{X}\mat{W}^{\top} + \vect{b}.
]

因此，阅读代码时不要只根据权重 tensor 的 shape 判断数学公式方向，而要结合框架 API 的约定。

\section{本书公式排版规范}

\subsection{公式与代码之间的对应关系}

本书尽量保持公式与代码之间的一一对应。例如，Scaled Dot-Product Attention 写作：
[
\Attention(\mat{Q}, \mat{K}, \mat{V})
=====================================

\softmax\left(
\frac{\mat{Q}\mat{K}^{\top}}{\sqrt{D_h}}
\right)\mat{V}.
]

对应到伪代码时，通常写为：

\begin{lstlisting}[language=Python,caption={Scaled Dot-Product Attention 的基本形式},label={lst:appendix-attention-basic}]
scores = q @ k.transpose(-2, -1)
scores = scores / math.sqrt(head_dim)
probs = torch.softmax(scores, dim=-1)
out = probs @ v
\end{lstlisting}

这里的 \texttt{transpose(-2, -1)} 对应公式中的 $\mat{K}^{\top}$，\texttt{dim=-1} 表示沿 key sequence 维度做 softmax。

\subsection{复杂度记号}

本书使用 $\mathcal{O}(\cdot)$ 表示渐进复杂度。若讨论 Transformer attention 的计算复杂度，常写作：
[
\mathcal{O}(B \cdot H \cdot S^2 \cdot D_h).
]

若讨论 KV Cache 显存复杂度，常写作：
[
\mathcal{O}(B \cdot L \cdot H_{\mathrm{kv}} \cdot S \cdot D_h).
]

需要注意，复杂度表达式只表达增长趋势，不等价于真实运行时间。真实性能还受到 kernel fusion、访存模式、HBM 带宽、Tensor Core 利用率、通信拓扑和调度策略影响。

\subsection{单位约定}

本书在系统性能章节中常用如下单位：

\begin{longtable}{p{0.22\textwidth}p{0.68\textwidth}}
\toprule
单位 & 含义 \
\midrule
FLOPs & 浮点运算次数。 \
TFLOPS & 每秒万亿次浮点运算。 \
GB / GiB & 十进制 GB 与二进制 GiB。工程排查中需注意二者差异。 \
GB/s & 每秒传输的字节数，常用于显存带宽、PCIe 带宽、网络带宽。 \
tokens/s & 推理或训练中每秒处理的 token 数。 \
steps/s & 训练中每秒完成的 optimizer step 数。 \
TTFT & Time To First Token，推理请求首 token 延迟。 \
TPOT & Time Per Output Token，生成阶段单 token 平均耗时。 \
\bottomrule
\end{longtable}

\section{附录 A 小结}

本附录定义了本书使用的主要数学符号、Transformer 维度命名、张量布局约定与公式排版规范。读者在阅读训练并行、推理优化、GPU kernel 和集群调度章节时，可以随时回到本附录对齐变量含义。尤其在分析 KV Cache、Tensor Parallelism、FlashAttention 与 Continuous Batching 时，明确维度语义往往比记住某个公式更重要。
